\documentclass[12pt]{article}
\usepackage[utf8]{inputenc}
\usepackage{geometry}
\geometry{a4paper, margin=1in}
\usepackage{enumitem}
\usepackage{titlesec}
\usepackage{xcolor}
\usepackage{amsmath}
\usepackage{amssymb}

% Formatting titles
\titleformat{\section}{\large\bfseries\color{blue!60!black}}{\thesection}{1em}{}
\titleformat{\subsection}{\normalsize\bfseries\color{blue!50!black}}{\thesubsection}{1em}{}

\title{\textbf{Reinforcement Learning (RL): Core Concepts}}
\author{}
\date{}

\begin{document}

\maketitle
\vspace{-1.5cm}

\section*{Reinforcement Learning (RL)}
\begin{itemize}[label=$\bullet$]
    \item The main concept on which RL works is an agent learning by trial and error.
    \item Unlike other types of machine learning, the agent is not explicitly told what the correct action is. It must discover which actions yield the most reward.
    \item Rewards may be delayed, meaning an action taken now can result in positive or negative rewards much later in the future (after multiple steps in a trajectory).
\end{itemize}

\section*{Environment \& The Loop}
\noindent\textbf{Environment:} The world that the agent lives in and interacts with. \\
\textbf{The Loop:} Agent sees a state $\rightarrow$ takes an action $\rightarrow$ receives a reward signal and the next state from the environment. \\
The environment can control when to ask the agent to act.

\section*{Key Terms}
\begin{enumerate}
    \item \textbf{State ($s$):} A complete description of the world.
    \item \textbf{Observation ($o$):} A partial description of the state.
    \item \textbf{Action Space:} The set of all valid actions available to the agent.
    \item \textbf{Trajectories:} A sequence of states and actions over time.
    \item \textbf{Policy ($\pi$):} Rule by which agent takes action.
    \item \textbf{Return ($G$):} Total reward over a trajectory.
    \item \textbf{Discount Factor ($\gamma$):} Factor to ensure the mathematical sum of future rewards converges.
\end{enumerate}

\section*{The Elements of Reinforcement Learning}

\subsection*{Policy}
\begin{itemize}[label=$\bullet$]
    \item A policy defines the learning agent's way of behaving at a given time. It is a mapping from the perceived states of the environment to the actions to be taken when in those states.
    \item A policy can be as simple as a basic lookup table or function, or it can involve extensive computation.
    \item The policy is the core of the RL agent because it alone is sufficient to determine behavior.
\end{itemize}

\subsection*{Return}
\begin{itemize}[label=$\bullet$]
    \item It is the Reward signal which defines the immediate goal of the RL problem.
    \item If an action chosen by the policy is followed by low reward, the policy will be updated to select a different action in that situation next time.
\end{itemize}

\subsection*{Value Functions}
Functions that estimate the expected return to predict how good a situation is:
\begin{itemize}[label=$\bullet$]
    \item \textbf{State-Value Function ($V$):} Measures how good it is to be in a specific state.
    \item \textbf{Action-Value Function ($Q$):} Measures how good it is to take a specific action in a specific state.
    \item \textbf{Optimal Value Functions ($V^*, Q^*$):} The maximum possible expected return an agent can achieve.
\end{itemize}

\subsection*{Exploration vs. Exploitation}
\begin{itemize}[label=$\bullet$]
    \item The fundamental trade-off the agent faces while learning.
    \item \textbf{Exploration:} Trying out new, untested actions to gather more information about the environment.
    \item \textbf{Exploitation:} Using known information to choose the action that is currently expected to yield the highest reward.
    \item An agent must balance both i.e. exploring too much wastes time, but exploiting too early means missing out on potentially better strategies.
\end{itemize}

\end{document}
